\section{Common coefficients and additional discriminants}
\label{sec:discriminants}

\subsection{A common CMB-to-galaxy normalization}
\label{sec:crossnorm}
The bare commitment ratio and galactic scale obey
\begin{equation}
\frac{\rho_{C,i}}{\rho_{b,i}}=\frac1\epsilon,\qquad
a_0(z)=\epsilon cH(z),\qquad
a_0(z)\frac{\rho_{C,i}}{\rho_{b,i}}=cH(z).
\end{equation}
The fitted formation convention includes the mass-conversion factor
$f_m=5.333\epsilon$, and release leaves the homogeneous survival factor
$s(z)=e^{-\eta_*F(z)}$. On the undepleted background branch, the
observable stock ratio therefore satisfies
\begin{equation}
a_0(z)\frac{\bar\rho_C(z)}{\bar\rho_b(z)}=f_m s(z)cH(z).
\label{eq:crossnorm}
\end{equation}
The fitted $H_0=68.9389\,\kms\,\mathrm{Mpc}^{-1}$ gives
$a_0(0)=1.2588\times10^{-10}\,\mathrm{m\,s^{-2}}$, about 2.3\%
above the reference value, without a separate fit of $a_0$.
This is a conditional relation between sectors. Cosmological abundance
and $H_0$ are inferred together; galaxy mass calibration limits the
elementary empirical comparison to roughly twenty percent. The phase
loading and mass convention retain their physical matching assumptions.
\source{U1, S02; T4--T5}

\subsection{Newton's constant and charged-lepton masses}
\label{sec:constants}
The selected particle-coefficient route gives
$G=6.6742890813\times10^{-11}\,\mathrm{m^3\,kg^{-1}\,s^{-2}}$,
$-0.073$ quoted CODATA standard deviations; the muon ratio is
$206.768280237$, at $-0.535$ quoted standard deviations. Its tau mass
is 1776.9188 MeV against $1776.93\pm0.09$ MeV. Of the 401 tested
routing integers, only the integer $7$ lies within two quoted standard deviations
for the stated route; of 43 tested Farey candidates, only $2/7$ lies
within one. These finite-menu uniqueness checks do not quantify the
historical search over possible formulas. The measured Newton scale
and marked-formula residuals were known during development.
\source{U1, S11}\cite{constantsnew}


\subsection{Halo tilt and a common Milky Way--LMC history}
\label{sec:halo}
Three eight-Gyr halo phase realizations retain late stock tilts of
$39.58^\circ$--$42.36^\circ$ and stellar tilts of
$21.38^\circ$--$22.33^\circ$. Rounding and precession persist.
These are repeatable outcomes of the specified kinetic model, with
force convergence and stationary formation still unresolved. The
aligned-branch failure remains part of the comparison. A common
Milky Way--LMC--stream fit remains open. The archived Galactic masses
inside 50 kpc exceed the cited constraint; Section~\ref{sec:jointtest}
quantifies that mismatch. \source{U1, S13; G3--G7}



\subsection{Old-disk vertical gravity}
\label{sec:olddisk}
At $a_0=1.2588\times10^{-10}\,\mathrm{m\,s^{-2}}$, the reduced
old-disk solver gives NGC~628 $\chi^2=2.04$ for five dispersions and
a $-0.21\sigma$ radial--vertical mass offset. The approximate algebraic
QUMOND comparison gives $\chi^2=8.95$ and $-2.91\sigma$. The small
normalization change preserves the earlier pattern. Gas-profile,
inclination, and scale-height uncertainties remain, and a round
conventional halo can give similar vertical gravity.
\source{U1, S15}


\subsection{Wide binaries and Solar-System null tests}
\label{sec:binaries}
With the parent-carrier rule, a wide binary has no independent capacity
halo. A radial Galactic-tide estimate for one to two solar masses at
10--30 kAU gives $\Delta v/v\simeq1.1\times10^{-5}$--$6\times10^{-4}$.
The range depends on separation, mass, and orientation; a universal
$10^{-4}$ ceiling is unwarranted. The selected Solar quadrupole
$Q_2=0$ is 0.89 quoted standard deviations from the Cassini central
value. No full Gaia selection-and-contamination likelihood has been
performed. \source{U1, S10}

